% frontmatter/reading-paths.tex
% Карта чтения: краткое описание пяти частей и логики переходов между ними.

\addchap{Как читать книгу}

Книгу можно читать линейно: каждая часть опирается на~предыдущие,
и~это самый надёжный маршрут. Если читать выборочно, главы первой
части обычно нужны как контекст для остальных~--- без них продвинутые
темы быстро становятся непрозрачными.

Часть~\ref{part:money-flows}, \enquote{Движение денег},
разбирает базовые понятия, без которых остальная книга превращается
в~набор разрозненных деталей: что считается деньгами в~банковской
инфраструктуре, кто кому должен и~в~какой момент, чем авторизация
отличается от~расчёта, как устроена сама карта и~её идентификаторы,
как выглядит жизненный цикл операции и~как кодируется протокол
ISO~8583. К~этим главам возвращаются все остальные части книги.

Часть~\ref{part:security-protocols}, \enquote{Безопасность
и~протоколы}, отвечает на~вопрос, почему платёжной системе вообще
стоит доверять. Карта доказывает подлинность чипом, токен заменяет
PAN, не~разрушая жизненный цикл операции, криптография даёт общий
язык секретов и~подписей, 3-D~Secure переносит проверку держателя
в~электронную коммерцию, а~PCI~DSS определяет цену хранения
чувствительных данных. Эти главы остаются на~уровне общих
механизмов; конкретные сети появятся в~следующей части.

Часть~\ref{part:payment-systems}, \enquote{Платёжные системы},
сравнивает эти механизмы в~конкретных системах:
Visa и~Mastercard, Мир, UnionPay, региональные карточные схемы СНГ,
СБП, мгновенные платежи мира, открытый банкинг, QR-кошельки, BNPL,
цифровые валюты центральных банков. К~каждой системе ставятся одни
и~те же вопросы: кто владеет правилами, какой документ главный,
как достигается окончательность расчёта, что меняется для банка,
продавца и~процессинга при переходе с~одной системы на~другую.
Глубина намеренно асимметрична: системы, которые чаще встречаются
русскоязычному инженеру, разобраны подробнее, чем международные
карточные схемы.

Часть~\ref{part:operations}, \enquote{Эксплуатация (ops)},
переходит с~языка \enquote{как должен работать платёж} на~язык того,
что происходит с~платежом уже в~проде. Регулирование, идентификация
клиентов и~ТСП, ПОД/ФТ, санкционный скрининг, фрод, подписки, сверка,
диспуты, управление отказами~--- темы, которые в~компании обычно
разнесены по~разным командам, но для платёжного продукта составляют
одну и~ту же реальность.

Часть~\ref{part:architecture}, \enquote{Архитектура}, складывает
компоненты обратно в~систему: платёжный шлюз, внутренний реестр
обязательств, мультивалютная логика, процессинговый центр. Это
не~рецепт \enquote{как собрать платформу с~нуля}, а~карта
ответственности~--- какой компонент за~что отвечает и~почему граница
между ними прошла именно здесь.

Если место в~продвинутой главе перестаёт быть понятным~--- скорее
всего, ответ находится в~одной из глав части~\ref{part:money-flows}.
Сначала~--- словарь движения денег, сообщений и~обязательств;
всё остальное собрано на~нём.
